% S3 Implementation Details
\subsection{Perception and Scene Memory}
\label{sec:supp-perception}
\noindent\textbf{Observation in simulation.} After every execution bundle and before every FM request, the
system renders segmentation masks from the scene cameras. An object is counted in a view if its mask there
covers at least 10 pixels (at least 50 in the grill scene), and $V_t$ contains the objects counted in at least
one view. The remaining facts come from the simulator: the region of each visible object is assigned by testing
its pose against the region bounding boxes, lid states are read from the lid joints, and placement-area overlap
is computed from the object's footprint when it was last visible. Gripper state is maintained by the executor.

\noindent\textbf{Identities and memory.} $K_t$ records \emph{which} objects have been observed, and $m_t$ stores
\emph{what} was last observed about each of them: its region $r_t(o)$ and the step $t$ of the observation.
After each observation,
\begin{equation}
m_t(o)=\begin{cases}
\big(r_t(o),\,t\big) & o\in V_t,\\
m_{t-1}(o) & o\in K_{t-1}\setminus V_t,
\end{cases}
\label{eq:supp-memory}
\end{equation}
so an occluded or removed-from-view object keeps its last region and observation time, and no entry is ever
deleted. The prompt lists such objects separately from the visible ones, with their last region and the number
of steps since they were seen. An object $o\notin K_t$ has no entry, never appears in a prompt, and a plan that
names it is rejected with the same message as for a non-existent name, so the planner cannot distinguish a
hidden object from an invalid one. A remembered object can be picked only if its last region is not closed off
by a closed lid. Without memory ($-$Memory), only $V_t$ is available.

\subsection{Replanning and Execution}
\label{sec:supp-replanning}
\noindent\textbf{\textit{IF}.} After each bundle, every $o\in K_t$ other than the lids and the held object is
assessed with its current region if visible and its remembered region otherwise. Let $C_d$ be the task-relevant
categories of the domain, $\mathit{cat}(o)$ the category of $o$ (mugs and groceries; raw and cooked meat), $r^\star(o)$ the goal region of $o$'s category
and $\mathcal{A}(\pi_t)$ the placement areas of the regions that some pending \texttt{place} in $\pi_t$ targets.
The trigger set is
\begin{equation}
\begin{aligned}
O_e=\big\{\,o \;\big|\;& \neg\,\mathit{cov}(o) \wedge \big[\big(\mathit{cat}(o)\in C_d \wedge \neg\,\mathit{done}(o)\big)\\
&\vee \big(\mathit{cat}(o)\notin C_d \wedge \mathit{fp}(o)\cap\mathcal{A}(\pi_t)\neq\emptyset\big)\big]\big\},
\end{aligned}
\label{eq:supp-trigger}
\end{equation}
where $\mathit{cov}(o)$ holds if a pending \texttt{pick} or \texttt{place} names $o$, or if the FM answered
\texttt{NO\_ACTIONS} for $o$ and its recorded region has not changed since; $\mathit{done}(o)$ holds if $o$'s
region is $r^\star(o)$ and, for meat, $h_t$ shows it cooked; and $\mathit{fp}(o)$ is its footprint. Coverage checks
only that a pending action names the object, not its destination or position. Obstruction is tested against
every pending placement regardless of its position in $\pi_t$, so the plate, which is not a goal category, triggers
only when it obstructs. Overcooking and serving raw are left to the evaluator and do not trigger. If
$O_e=\emptyset$, execution continues without an FM call.

\noindent\textbf{\textit{WHEN}.} A trigger request is sent to the FM while the robot executes the bundles of
$\pi_t$ that cannot interfere with the correction. The \emph{affected set} contains the objects of $O_e$, their
goal regions, the staging region $S$ (the area where a correction may set objects aside), the container of a
trigger object whose placement area it overlaps, and the lids that close off an affected region or a trigger
object's container (closing it would change the object's state, e.g.\ cook a cooked meat again). For bundle
$b$, let $\mathit{obj}(b)$ be its objects, $\mathit{src}(b)$ the region its pick takes from (projected along the
plan, so a second pick of the same object uses the region the first placement put it in) and $\mathit{dst}(b)$
its placement region or, for a lid action, the regions the lid closes off. Bundles are scanned in plan order
with blocked sets initialized to $B_o$ = affected objects, $B_d$ = affected regions and $B_s=B_d\setminus S$:
\begin{equation}
b \text{ eligible} \iff \mathit{obj}(b)\cap B_o=\emptyset \;\wedge\; \mathit{dst}(b)\cap B_d=\emptyset
\;\wedge\; \mathit{src}(b)\cap B_s=\emptyset .
\label{eq:supp-eligible}
\end{equation}
A withheld bundle propagates its dependencies to later bundles:
$B_o\!\leftarrow\!B_o\cup\mathit{obj}(b)$, $B_d\!\leftarrow\!B_d\cup\mathit{dst}(b)\cup\mathit{src}(b)$ and
$B_s\!\leftarrow\!B_s\cup\mathit{dst}(b)$. Later bundles that use its objects, place into or take from its
destination, or place into the region it would vacate are therefore also withheld. Because $S\in B_d$ but
$S\notin B_s$, placements into the staging region wait, while removing an object from it remains allowed.
Eligible bundles run in their original relative order and may precede an earlier withheld bundle. The FM request
lists them as scheduled and excludes them from the remaining plan shown, and the scene description is not
advanced by their effects. Execution checks for the response only between bundles; if no eligible bundle is
left, the robot waits. Trigger checks are suspended during the wait and resume after the merge. Without
concurrency ($-$WHEN), execution stops during inference.

\noindent\textbf{\textit{WHERE}.} A trigger request returns corrective blocks for the objects of $O_e$, not a
new task plan. Each block lists the objects it handles, an urgency, an insertion position (front, after a named
pending action, or end), a one-sentence reason and its actions; a block with \texttt{NO\_ACTIONS} states that
its objects need none. Let $\pi_{t'}$ be the remaining plan after the actions completed during inference are
removed. Only actions of $\pi_{t'}$ that were not scheduled for concurrent execution are eligible anchors; an
anchor on a scheduled action is rejected even if that action has not yet run. With urgent blocks
$U_1,\dots,U_k$ in the order returned and deferred blocks $D_j$ placed directly after their anchor or appended,
\begin{equation}
\widehat{\pi}=U_1\cdots U_k\;\cdot\;\mathrm{ins}\big(\pi_{t'},\{D_j\}\big),
\qquad \pi_{t'}\ \text{is a subsequence of}\ \widehat{\pi},
\label{eq:supp-merge}
\end{equation}
so corrective actions are added while the pending actions keep their relative order and identifiers. The
planner's urgency and position are used as returned; urgent blocks must be placed at the front.

\noindent\textbf{Validation and recovery.} A correction is rejected if its blocks are malformed, handle an
object outside $O_e$, leave an object of $O_e$ unhandled, act on objects other than those they list, contain an
incomplete pick--place sequence, or anchor on an ineligible action; if the merged plan fails the plan check
(unknown or never-observed names, inconsistent pick--place order, an inaccessible remembered object); or if it
places into a region before an overlapping trigger object has been removed from that region's placement area
(insertion too late, using the overlap recorded at the trigger). No check tests whether a correction is
inserted too early or whether its urgency respects a procedure such as cooking; those remain FM decisions. A
rejected correction is re-queried as the same corrective request with the failure code and the offending
element added, and the remaining plan is kept. An execution failure (a failed grasp, placement or motion plan,
or a violated precondition such as a closed lid) instead requests a full replacement of the remaining plan from
the updated scene and history; a failure of a concurrently executed bundle discards the pending correction and
is handled the same way. There are no task-level local retries. An inconsistent executor state, such as placing
an object that is not held, ends the trial.

\noindent\textbf{Termination.} At most $N=11$ FM calls are made per trial (the initial request and ten
replanning calls, re-queries included). A trial also ends when the plan is exhausted with no pending trigger,
when a full-plan request returns \texttt{NO\_ACTIONS}, on a terminal failure, or when the planner repeats an
earlier answer for the same state and failure a second time (the first repetition is re-queried with a note). Termination is not
success: success is evaluated afterwards from the final state and execution history with $G_\tau$, and no
runtime goal check was used in the reported runs.

\subsection{Simulation Controllers}
\label{sec:supp-controllers}
A Franka Panda executes the primitives in MuJoCo. Inverse kinematics uses batched damped least squares and
collision-free paths use RRT-Connect. In the kitchen, each pick--place bundle is planned with PDDLStream (grasp,
inverse-kinematics and motion streams; at most 60\,s per solve), with placement poses sampled inside the
target region's placement area; the mug stored in the cupboard is picked with a scripted motion, and the box lid
is opened by replaying a prerecorded trajectory. In the grill, transfers are executed in stages that are planned
as the robot moves, every meat is placed at the fixed pose of the grill placement area, and the grill-lid and
plate motions replay prerecorded trajectories. Pre- and post-action checks classify failures such as a closed lid
at a placement, a failed grasp or a placement outside the target region.

\subsection{Real-Robot Execution}
\label{sec:supp-real}
The physical system uses an xArm manipulator and three Intel RealSense D435 cameras. SAM2.1 provides object
segmentation masks, observed objects are associated with predefined placement regions, and these object--region
observations update the same memory layer as in simulation. Point clouds from the cameras support grasp
generation.

\medskip
\noindent\fbox{\parbox{\dimexpr\columnwidth-2\fboxsep-2\fboxrule}{\todo{Real-robot details to be added:
the five settings, region definitions, perception and grasping pipeline, and the trials performed.}}}
